\section{Secondary quantitative routes and local stability}
\label{sec:two-uniform}

The rigidity theorem uses the stabilizer hypothesis to identify the
symmetry group.  Finiteness of that group needs far less: two-party
marginals alone, at any local dimension, with no stabilizer structure.
Throughout this subsection \(q\geq2\) is an arbitrary integer; only the
statements naming Clifford operators return to \(q=p^e\).

Call a pure state
\(\lvert\psi\rangle\in(\mathbb C^q)^{\otimes n}\) \emph{\(2\)-uniform} if
\[
 \rho^\psi_{\{j,k\}}=q^{-2}I\qquad\text{for all }j\ne k .
\]
Tracing once more makes every one-party marginal \(q^{-1}I\).  Such states
exist only for \(n\geq4\), since at \(n=3\) purity forces
\(\rank\rho_{\{1,2\}}=\rank\rho_{\{3\}}\leq q\); and every
\(\AME(2m,q)\) state with \(m\geq2\) is \(2\)-uniform, stabilizer or not.
Write
\[
 G(\psi)=\{U_1\otimes\cdots\otimes U_n:
   (U_1\otimes\cdots\otimes U_n)\lvert\psi\rangle
     \in\mathbb C\lvert\psi\rangle\}
\]
for its product-unitary symmetry group and, for a product unitary \(U\),
\[
 \varepsilon(U)=\min_\theta\bigl\lVert U\lvert\psi\rangle
    -e^{i\theta}\lvert\psi\rangle\bigr\rVert
  =\bigl(2-2\lvert\langle\psi\rvert U\lvert\psi\rangle\rvert
   \bigr)^{1/2}
\]
for its defect.  Because operators on distinct sites commute, a product
unitary is a single exponential:
\(\bigotimes_je^{iH_j}=e^{iL}\) with \(L=\sum_jH_j^{(j)}\), where
\(H^{(j)}\) denotes \(H\) acting at site \(j\).  Splitting each \(H_j\)
into a traceless part \(h_j\) and a multiple of the identity leaves
\(L=M+cI\) with \(M=\sum_jh_j^{(j)}\).  Everything below rests on one
computation about \(M\).

\AMEFigureFour

\begin{lemma}[Local-generator isometry]
\label{lem:local-generator-isometry}
Let \(\lvert\psi\rangle\) be \(2\)-uniform and let \(h_j\) and \(h'_j\),
\(1\leq j\leq n\), be traceless Hermitian operators on \(\mathbb C^q\).
Put \(M=\sum_jh_j^{(j)}\) and \(M'=\sum_jh_j'^{(j)}\).  Then
\(\langle\psi\rvert M\lvert\psi\rangle=0\) and
\[
 \bigl\langle M\lvert\psi\rangle,\,M'\lvert\psi\rangle\bigr\rangle
  =\frac1q\sum_{j=1}^n\Tr(h_jh'_j).                      \tag{B.1}
\]
In particular \(\lVert M\lvert\psi\rangle\rVert^2
 =q^{-1}\sum_j\lVert h_j\rVert_F^2\), so
\((h_1,\ldots,h_n)\mapsto\sqrt q\,M\lvert\psi\rangle\) is an isometry from
the traceless local generators, carrying the product Frobenius form, into
the Hilbert space regarded as a real inner product space.
\end{lemma}

\begin{proof}
The one-party marginals give
\(\langle M\rangle=\sum_j\Tr(\rho^\psi_{\{j\}}h_j)
 =\sum_j\Tr(h_j)/q=0\).  Expand \(\langle\psi\rvert MM'\lvert\psi\rangle\).
The terms with equal site indices contribute
\(\Tr(\rho^\psi_{\{j\}}h_jh'_j)=\Tr(h_jh'_j)/q\).  A term with \(j\ne k\)
contributes
\(\Tr\bigl(\rho^\psi_{\{j,k\}}(h_j\otimes h'_k)\bigr)
 =\Tr(h_j)\Tr(h'_k)/q^2=0\); both maximal mixedness of the pair marginal
and tracelessness are used.  Hence
\(\langle\psi\rvert MM'\lvert\psi\rangle=q^{-1}\sum_j\Tr(h_jh'_j)\), which
is real.  Since \(M\) and \(M'\) are Hermitian, the imaginary part of
\(\langle M\psi,M'\psi\rangle\) is
\(\langle[M,M']\rangle/2i\), and it vanishes as well: operators at distinct
sites commute, and the surviving one-site commutators are traceless, so
each term is of the form \(\Tr(\rho^\psi_{\{j\}}[h_j,h'_j])\)
\(=\Tr([h_j,h'_j])/q=0\).  A real symmetric bilinear form agreeing with
the product Frobenius form is that form.
\end{proof}

The identity has no error term, and that is the content of the
\(2\)-uniformity hypothesis rather than a coincidence: the cross terms
cancel exactly because the pair marginals are exactly \(q^{-2}I\).  A Bell
pair is the boundary case.  There \(\rho_{\{1,2\}}\) is a pure state, the
cross term survives, and \(M=h\otimes I-I\otimes h^{T}\) annihilates the
state for every \(h\) --- the \(U\otimes\overline U\) gauge line that makes
\(m\geq2\) sharp in Theorem~\ref{thm:lu-lc-rigidity}.

The second ingredient is that a one-parameter group of product unitaries
has a generator of summed local form.  This is not automatic: a
one-parameter subgroup of \(G(\psi)\) is a priori only a continuous
homomorphism into \(U(q^n)\), and the factors of a product unitary are
determined only up to phases, so differentiating factor by factor carries a
lifting obligation.

\begin{lemma}[Product generators]\label{lem:product-lie}
Let \(\pi:U(q)^n\to U(q^n)\) send \((U_1,\ldots,U_n)\) to
\(U_1\otimes\cdots\otimes U_n\).  Then \(\pi\) is a smooth homomorphism of
compact Lie groups, its image \(P\) is a closed subgroup of \(U(q^n)\), and
\[
 \operatorname{Lie}(P)
  =\Bigl\{i\textstyle\sum_{j}H_j^{(j)}:
    H_j=H_j^\dagger\ \text{on }\mathbb C^q\Bigr\}.       \tag{B.2}
\]
Hence every one-parameter subgroup of a closed subgroup of \(P\) has a
generator of that form.
\end{lemma}

\begin{proof}
The tensor product of unitaries is multiplicative factor by factor, so
\(\pi\) is a smooth homomorphism.  Its source is compact, so \(P\) is
compact and therefore closed, and Cartan's theorem makes \(P\) an embedded
Lie subgroup of \(U(q^n)\).  A surjective homomorphism of Lie groups has
surjective differential, so
\(\operatorname{Lie}(P)=d\pi\bigl(\mathfrak u(q)^n\bigr)\), which is the
displayed space; its kernel is
\(\{(i\tau_1I,\ldots,i\tau_nI):\sum_j\tau_j=0\}\).  A closed subgroup
\(H\leq P\) is a Lie subgroup with
\(\operatorname{Lie}(H)\subseteq\operatorname{Lie}(P)\), and the generator
of a one-parameter subgroup of \(H\) lies in \(\operatorname{Lie}(H)\).
\end{proof}

\begin{theorem}[Discreteness from \(2\)-uniformity]
\label{thm:two-uniform-discrete}
Let \(\lvert\psi\rangle\) be a \(2\)-uniform pure state of \(n\) qudits of
local dimension \(q\).  Then \(\operatorname{Lie}(G(\psi))=i\mathbb RI\),
and \(G(\psi)/U(1)\cdot I\) is finite.  In particular every
\(\AME(2m,q)\) state with \(m\geq2\), stabilizer or not, has finite
product-unitary symmetry group modulo global phase.
\end{theorem}

\begin{proof}
The set \(\{U\in U(q^n):U\lvert\psi\rangle\in\mathbb C\lvert\psi\rangle\}\)
is a closed subgroup, so \(G(\psi)=P\cap\{\cdots\}\) is a closed subgroup
of \(P\) and hence a compact Lie group.  Let
\(X\in\operatorname{Lie}(G(\psi))\).  By Lemma~\ref{lem:product-lie},
\(X=iL\) with \(L=\sum_jH_j^{(j)}\) Hermitian; write \(L=M+cI\) as above.
Each \(e^{itL}\) fixes the ray of \(\lvert\psi\rangle\), so
\(e^{itL}\lvert\psi\rangle=\lambda(t)\lvert\psi\rangle\) with
\(\lambda(t)=\langle\psi\rvert e^{itL}\lvert\psi\rangle\) unimodular and
smooth in \(t\).  Differentiating at \(t=0\) shows that
\(\lvert\psi\rangle\) is an eigenvector of the Hermitian operator \(L\),
say \(L\lvert\psi\rangle=\mu\lvert\psi\rangle\) with \(\mu\) real.  Pairing
with \(\lvert\psi\rangle\) and using \(\langle M\rangle=0\) gives
\(\mu=c\), hence \(M\lvert\psi\rangle=0\).  By
Lemma~\ref{lem:local-generator-isometry}, \(\sum_j\lVert h_j\rVert_F^2=0\),
so every \(h_j=0\) and every \(H_j\) is scalar.  Thus
\(\operatorname{Lie}(G(\psi))=i\mathbb RI\).  Since \(U(1)\cdot I\) is a
closed central one-dimensional subgroup of \(G(\psi)\), the compact Lie
group \(G(\psi)/U(1)\cdot I\) has dimension zero and is therefore finite.
\end{proof}

\begin{remark}[Which torus is divided out]
\label{rem:local-phase-torus}
Upstairs in \(U(q)^n\) the preimage \(\pi^{-1}(G(\psi))\) always contains
the local-phase torus \(T=U(1)^n\), of dimension \(n\).  This does not
contradict the theorem: \(\pi(T)=U(1)\cdot I\) and \(\ker\pi\subset T\) has
dimension \(n-1\), so
\(\pi^{-1}(G(\psi))/T\cong G(\psi)/U(1)\cdot I\).  Statements about local
factors below are statements on \(\mathrm{PU}(q)^n\), not statements
modulo one global phase.
\end{remark}

\begin{remark}[Discreteness in the literature]
\label{rem:discreteness-prior-art}
Finiteness of this kind is known in two special cases.  Wirthm\"uller
determines the identity component of the local-unitary automorphism group
of a binary stabilizer code and gives a protected-qubit criterion for
finiteness of the induced group
\cite[Theorem~7 and Corollary~11]{Wirthmuller2011}; his exclusion of Bell
factors is the \(n\geq4\) boundary above.  Englbrecht and Kraus determine
the local symmetries of qubit stabilizer states, including where continuous
ones occur \cite{EnglbrechtKraus2020}.  Tan computes the local symmetry
group of the four-qutrit AME state outright, of order \(5832\)
\cite[Theorem~B.3]{Tan2026}.  Theorem~\ref{thm:two-uniform-discrete} covers
arbitrary local dimension and arbitrary, in particular nonstabilizer,
\(2\)-uniform states; to our knowledge no statement at that scope is
available.

The hypothesis is exactly \(2\)-uniformity, and one-party marginals are
not enough.  S{\l}owik, Sawicki, and Maci\k{a}\.zek construct locally
maximally entangled states --- every one-party marginal maximally mixed
--- whose local-unitary symmetry group contains a nontrivial continuous
image of \(\mathrm{SU}(N)\) and is therefore positive-dimensional, with
dimension growing in the number of parties
\cite[Lemma~2]{SlowikSawickiMaciazek2021}.  Passing from one-party to
two-party marginals is what forces discreteness.  In the language of
geometric invariant theory the theorem says that a \(2\)-uniform state is a
stable point for \(\mathrm{SL}(q,\mathbb C)^n\) rather than merely a
critical one; stabilizer dimensions there are known for generic states and
in the extremal qubit case
\cite{BryanReichsteinVanRaamsdonk2018,%
BryanLeutheusserReichsteinVanRaamsdonk2019,LyonsWalckBlanda2008,%
GourKrausWallach2017}, but a genericity statement does not decide a
particular state.
\end{remark}

\begin{remark}[Two-uniformity does not identify the group]
\label{rem:two-uniform-not-clifford}
An explicit state shows that the second half of the rigidity phenomenon
needs strictly more than the first.  Let \(q=2\) and
\(C=\mathrm{RM}(1,4)\leq\F_2^{16}\), whose nonzero codeword weights are
\(8\) and \(16\) and whose dual \(\mathrm{RM}(2,4)\) has minimum distance
\(4\).  The equal-phase CSS state of \(C\) has label group
\(C_X\oplus C^\perp_Z\), so by \((2.2)\) a marginal fails to be maximally
mixed only when a nonzero word of \(C\) or of \(C^\perp\) is supported
inside it; both distances exceed \(2\), so the state is \(2\)-uniform and
Theorem~\ref{thm:two-uniform-discrete} makes its symmetry group finite.
That group is not Clifford.  Writing \(T=\diag(1,e^{i\pi/4})\), the product
\(T^{\otimes16}\) multiplies \(\lvert x\rangle\) by
\(e^{i\pi\operatorname{wt}(x)/4}\), which is \(1\) at every codeword
weight, so \(T^{\otimes16}\) fixes the state exactly; and
\(TX(1)T^\dagger=2^{-1/2}\bigl(X(1)+iX(1)Z(1)\bigr)\) is not a Weyl axis.
Two-uniformity buys finiteness and nothing else.  Identifying the finite
group as local Clifford is what the AME hypothesis and
Theorem~\ref{thm:lu-lc-rigidity} supply, and the next corollary is that
identification.
\end{remark}

\begin{corollary}[Discrete product-unitary symmetry]
\label{cor:discrete-lu-symmetry}
Fix a stabilizer AME state \(\lvert\psi\rangle\) as in
Theorem~\ref{thm:lu-lc-rigidity}, and define the local-factor tuple groups
\[
\begin{split}
G&=\{(U_i)\in U(q)^{2m}:
       (\bigotimes_iU_i)\lvert\psi\rangle
          \in\mathbb C\lvert\psi\rangle\},\\
\widetilde G
 &=\{((U_i),\pi)\in U(q)^{2m}\rtimes S_{2m}:
       (\bigotimes_iU_i)\pi\lvert\psi\rangle
          \in\mathbb C\lvert\psi\rangle\},
\end{split}
\]
where \(S_{2m}\) permutes the local factors.  Let
\(T=(S^1)^{2m}\) be the coordinatewise scalar subgroup and put
\(\Gamma=G/T\), \(\widetilde\Gamma=\widetilde G/T\).  Then
\[
  1\longrightarrow T\longrightarrow G\longrightarrow \Gamma
    \longrightarrow 1,
  \qquad
  1\longrightarrow T\longrightarrow \widetilde G
    \longrightarrow \widetilde\Gamma\longrightarrow 1
\]
are short exact sequences with closed connected kernel and finite
discrete quotient; both inclusions of \(T\) are continuous.  In
particular, \(G\) and \(\widetilde G\) are finite unions of connected
components, each a translate of \(T\).  Thus the local-factor
automorphism group modulo one-site scalar phases is finite and discrete,
and its identity component before quotienting is exactly \(T\).
If \(\Pi\) is the subgroup of
party permutations realized by \(\widetilde G\), there is a further exact
sequence
\[
  1\longrightarrow \Gamma\longrightarrow\widetilde\Gamma
    \longrightarrow\Pi\longrightarrow 1.
\]
This last extension splits exactly when the realized permutations admit a
homomorphic choice of projective product-unitary lifts.  We use only the
finiteness statement here; the extension class is not part of the
quantitative argument.
\end{corollary}

\begin{proof}
Every local factor in a tuple from \(G\) or \(\widetilde G\) is
Clifford.  The projective
single-qudit Clifford group over a finite field is finite, as is the
party-permutation group.  The kernel of projectivizing each local factor
is exactly the torus of one-site scalar phases.  That torus is the
identity component, hence is closed; the finite Hausdorff quotient is
discrete.  Translation identifies every connected component with the
scalar torus, so finitely many quotient representatives cover the group.
Forgetting the local factors maps
\(\widetilde\Gamma\) onto the realized permutation subgroup \(\Pi\).
Its kernel consists exactly of the fixed-party projective
automorphisms, which gives the third sequence.  A splitting is, by
definition, a group-homomorphic right inverse; arbitrary or
generator-by-generator lifts do not provide one.
\end{proof}


Discreteness makes the exact symmetries isolated.  The same identity
\((B.1)\) says how fast the defect grows away from them, with a constant
that does not involve the number of parties.

\begin{theorem}[Local stability]\label{thm:two-uniform-stability}
Let \(\lvert\psi\rangle\) be \(2\)-uniform, let each \(h_j\) be traceless
Hermitian on \(\mathbb C^q\), and put
\[
 U=\bigotimes_{j=1}^ne^{ih_j},\qquad
 t=\sum_{j=1}^n\lVert h_j\rVert_{\mathrm{op}},\qquad
 D=\Bigl(\sum_{j=1}^n\lVert h_j\rVert_F^2\Bigr)^{1/2}.
\]
If \(t\leq\frac12\) then
\[
 \frac{D^2}q\Bigl(1-\frac t3\Bigr)\leq\varepsilon(U)^2
 \leq\frac{D^2}q\Bigl(1+\frac t3\Bigr).                  \tag{B.3}
\]
In particular \(D\leq\sqrt{6q/5}\,\varepsilon(U)<1.096\sqrt q\,
\varepsilon(U)\), and \(\varepsilon(U)^2/D^2\to q^{-1}\) as \(t\to0\),
uniformly in \(n\).
\end{theorem}

\begin{proof}
Since the sites commute, \(U=e^{iM}\) with \(M=\sum_jh_j^{(j)}\), and
\(\lVert M\rVert_{\mathrm{op}}\leq t\) by the triangle inequality.  Write
\(\langle\cdot\rangle=\langle\psi\rvert\cdot\lvert\psi\rangle\).  Taylor's
formula with integral remainder gives
\(e^{ix}=1+ix-x^2/2+r(x)\) with \(\lvert r(x)\rvert\leq\lvert x\rvert^3/6\)
for real \(x\).  Let \(M=\sum_i\lambda_iP_i\) be the spectral decomposition
and put \(p_i=\langle P_i\rangle\geq0\), so that \(\sum_ip_i=1\).  Then
\(\langle r(M)\rangle=\sum_ip_ir(\lambda_i)\) and
\[
 \lvert\langle r(M)\rangle\rvert
  \leq\frac16\sum_ip_i\lvert\lambda_i\rvert^3
  \leq\frac{t\langle M^2\rangle}6 .
\]
By Lemma~\ref{lem:local-generator-isometry}, \(\langle M\rangle=0\) and
\(\langle M^2\rangle=D^2/q\), so
\(\langle e^{iM}\rangle=1-\langle M^2\rangle/2+\langle r(M)\rangle\).
Also \(\langle M^2\rangle\leq t^2\leq\frac14\), so
\(1-\langle M^2\rangle/2>0\) and
\[
 \Bigl\lvert\,\lvert\langle e^{iM}\rangle\rvert
   -\bigl(1-\tfrac12\langle M^2\rangle\bigr)\Bigr\rvert
 \leq\frac{t\langle M^2\rangle}6 ,
\]
using \(\lvert\langle e^{iM}\rangle\rvert\geq
\operatorname{Re}\langle e^{iM}\rangle\) for the lower estimate.  Since
\(\varepsilon(U)^2=2-2\lvert\langle e^{iM}\rangle\rvert\), this is exactly
\((B.3)\).  At \(t\leq\frac12\) the left inequality reads
\(\varepsilon^2\geq\frac56D^2/q\), and the two-sided form gives the
limit.
\end{proof}

The constant \(\sqrt{6/5}=1.0954\ldots\) is therefore sharp to within about
ten percent: the true asymptotic constant is \(\sqrt q\), and the excess is
the cubic Taylor cutoff, not structure.

\begin{remark}[The stability constant is a Fisher metric]
\label{rem:fisher-isotropy}
Identity \((B.1)\) also explains why the constant does not involve \(n\).
For the family \(U(s)=\bigotimes_je^{ish_j}=e^{isM}\), the quantum Fisher
information of a pure state is \(4\operatorname{Var}_\psi(M)\)
\cite{BraunsteinCaves1994}, so
\[
 F_Q=\frac4q\sum_{j=1}^n\lVert h_j\rVert_F^2 .
\]
Splitting the variance of a sum of one-site terms into single-party
variances and cross-party covariances is the standard framework for
metrologically useful entanglement
\cite{HyllusEtAl2012,Toth2012}; \(2\)-uniformity is exactly the hypothesis
that annihilates every covariance and pins every single-party variance.  On
traceless local generators --- and only there, since \((B.1)\) says
nothing about generators that are not sums of one-site terms --- the Fisher
form is a flat multiple of the Euclidean one.  A degenerate direction would
be a continuous product symmetry, which
Theorem~\ref{thm:two-uniform-discrete} excludes; an enhanced direction, of
the kind a GHZ state supplies, would need the pair correlations that
\(2\)-uniformity forbids.  The stability constant is the inverse of that
flat metric.

The Fisher form is quadratic data, so it governs the constant and says
nothing about how large a neighbourhood the quadratic estimate covers.
That size is fixed by the higher moments of \(M\), and
Theorem~\ref{thm:k-uniform-stability} and
Proposition~\ref{prop:stability-region} below determine it.
\end{remark}

The size of the certified region is controlled by the state's uniformity
order, not by the number of parties.  Call
\(\lvert\psi\rangle\) \emph{\(k\)-uniform} if every reduced operator on at
most \(k\) parties is maximally mixed; \(2\)-uniformity is the case
\(k=2\), and an \(\AME(2m,q)\) state is \(m\)-uniform.  The hypothesis
\(t\leq\frac12\) in Theorem~\ref{thm:two-uniform-stability} pays for a
cubic Taylor remainder.  Each additional order of uniformity removes one
order from that remainder, because the state and the tracial state assign
the same value to every power of \(M\) below the uniformity order.

\begin{lemma}[Tracial agreement below the uniformity order]
\label{lem:tracial-agreement}
Let \(\lvert\psi\rangle\) be \(k\)-uniform on \(n\) qudits.  Then
\(\langle\psi\rvert O\lvert\psi\rangle=q^{-n}\Tr O\) for every operator
\(O\) supported on at most \(k\) parties.  In particular, with
\(M=\sum_jh_j^{(j)}\) as above,
\(\langle\psi\rvert M^r\lvert\psi\rangle=q^{-n}\Tr(M^r)\) for every
\(r\leq k\).
\end{lemma}

\begin{proof}
If \(O=O_S\otimes I\) with \(\lvert S\rvert\leq k\) then
\(\langle O\rangle=\Tr(\rho^\psi_SO_S)=q^{-\lvert S\rvert}\Tr O_S
=q^{-n}\Tr O\).  Expanding \(M^r\) gives a sum of products of one-site
operators touching at most \(r\leq k\) sites, so every term is of that
shape.
\end{proof}

The tracial state \(q^{-n}\Tr\) is the product of the one-site tracial
states, and the site summands of \(M\) commute, so
\(q^{-n}\Tr f(M)=\mathbb E\,f(X_1+\cdots+X_n)\) with the \(X_j\)
independent and \(X_j\) uniform on the spectrum of \(h_j\).  The tracial
main term is therefore a product of one-site characteristic functions, and
the following contraction estimate controls each factor.

\begin{lemma}[Per-site contraction]\label{lem:site-contraction}
Let \(h\) be traceless Hermitian on \(\mathbb C^q\) with spectral spread
\(s=\lambda_{\max}(h)-\lambda_{\min}(h)\).  Then
\[
 \bigl\lvert q^{-1}\Tr e^{ih}\bigr\rvert^2
   \leq1-c(s)\lVert h\rVert_F^2/q,
 \qquad
 c(s)=\max\Bigl\{1-\tfrac{s^2}{12},\ \tfrac4{\pi^2}\,[s\leq\pi]\Bigr\}.
\]
\end{lemma}

\begin{proof}
Write \(q^{-1}\Tr e^{ih}=\mathbb E\,e^{iX}\) with \(X\) uniform on the
spectrum.  For an independent copy \(X'\),
\(\lvert\mathbb E\,e^{iX}\rvert^2=\mathbb E\cos(X-X')\) and
\(\lvert X-X'\rvert\leq s\).  From
\(\cos u\leq1-u^2/2+u^4/24=1-(u^2/2)(1-u^2/12)\) and \(u^2\leq s^2\),
\(\mathbb E\cos(X-X')\leq1-(1-s^2/12)\operatorname{Var}X\); from
\(\cos u\leq1-2u^2/\pi^2\) for \(\lvert u\rvert\leq\pi\),
\(\mathbb E\cos(X-X')\leq1-(4/\pi^2)\operatorname{Var}X\).  Tracelessness
gives \(\operatorname{Var}X=\Tr(h^2)/q=\lVert h\rVert_F^2/q\).
\end{proof}

\begin{theorem}[Stability at uniformity order \(k\)]
\label{thm:k-uniform-stability}
Let \(\lvert\psi\rangle\) be \(k\)-uniform with \(k\geq2\), let each
\(h_j\) be traceless Hermitian with spectral spread at most \(s\), and keep
\(t\), \(D\), \(U\) as in Theorem~\ref{thm:two-uniform-stability}.  Put
\[
 c=c(s),\qquad
 \lambda=\frac{2t^{\,k-1}}{(k+1)!},\qquad
 \rho=c-\frac{c^2D^2}{4q}-2\lambda .
\]
Then \(\varepsilon(U)^2\geq\rho D^2/q\); in particular \(\rho>0\) gives
\(D\leq\sqrt{q/\rho}\;\varepsilon(U)\).
\end{theorem}

\begin{proof}
Write \(e^{ix}=\sum_{r\leq k}(ix)^r/r!+\varrho(x)\) with
\(\lvert\varrho(x)\rvert\leq\lvert x\rvert^{k+1}/(k+1)!\).  By
Lemma~\ref{lem:tracial-agreement} the state and the tracial state agree on
every \(M^r\) with \(r\leq k\), so
\(\langle e^{iM}\rangle-q^{-n}\Tr e^{iM}
=\langle\varrho(M)\rangle-q^{-n}\Tr\varrho(M)\).  For any spectral weight
vector \(p_i\geq0\) with \(\sum_ip_i=1\),
\[
 \Bigl\lvert\sum_ip_i\varrho(\lambda_i)\Bigr\rvert
 \leq\frac{\lVert M\rVert_{\mathrm{op}}^{\,k-1}}{(k+1)!}\sum_ip_i\lambda_i^2 ,
\]
and both states give \(M^2\) the value \(D^2/q\), by
Lemma~\ref{lem:tracial-agreement} with \(r=2\) on one side and by
tracelessness on the other.  Since
\(\lVert M\rVert_{\mathrm{op}}\leq t\), each term is at most
\(t^{k-1}D^2/(q(k+1)!)\), so the two expectations differ by at most
\(\lambda D^2/q\).  The sites commute, so
\(q^{-n}\Tr e^{iM}=\prod_jq^{-1}\Tr e^{ih_j}\), and
Lemma~\ref{lem:site-contraction} with \(1-y\leq e^{-y}\) gives
\[
 \bigl\lvert\langle e^{iM}\rangle\bigr\rvert
  \leq\prod_j\bigl(1-c\lVert h_j\rVert_F^2/q\bigr)^{1/2}
      +\lambda D^2/q
  \leq e^{-cD^2/(2q)}+\lambda D^2/q .
\]
With \(x=cD^2/(2q)\) and \(1-e^{-x}\geq x-x^2/2\),
\(\varepsilon(U)^2=2-2\lvert\langle e^{iM}\rangle\rvert
\geq2x-x^2-2\lambda D^2/q=\rho D^2/q\).
\end{proof}

\begin{corollary}[The certified radius grows with the uniformity order]
\label{cor:k-uniform-region}
Let \(\lvert\psi\rangle\) be \(k\)-uniform with \(k\geq2\).  If every
\(h_j\) has spectral spread at most \(\frac12\), if \(D\leq\sqrt q/2\), and
if
\[
 t\leq R_k:=\Bigl(\frac{(k+1)!}{48}\Bigr)^{1/(k-1)},
\]
then \(D\leq\sqrt{6q/5}\,\varepsilon(U)\).  One has
\(\log R_k=\log k-1+O(\log k/k)\), so \(R_k\sim k/e\).
\end{corollary}

\begin{proof}
Spread at most \(\frac12\) gives \(c\geq1-\frac1{48}=0.97917\); with
\(D^2\leq q/4\), \(c^2D^2/(4q)\leq c^2/16\), so
\(c-c^2D^2/(4q)\geq0.91924\); and \(t^{k-1}\leq(k+1)!/48\) gives
\(2\lambda\leq\frac1{12}\).  Hence \(\rho\geq0.83591>\frac56\) and
\(\sqrt{q/\rho}\leq\sqrt{6q/5}\).  The asymptotics are Stirling's formula
applied to \(\log R_k=(\log(k+1)!-\log48)/(k-1)\).
\end{proof}

At \(k=2\) this is weaker than Theorem~\ref{thm:two-uniform-stability},
which reaches \(t\leq\frac12\) while \(R_2=\frac18\); the loss is the
passage through the tracial main term, which the direct argument avoids.
The gain begins at \(k=3\), where \(R_3=1/\sqrt2\), and grows from there:
\(R_4=1.357\), \(R_{10}=4.548\), \(R_{30}=12.93\).  For a stabilizer
\(\AME(2m,q)\) state \(k=m\), so the certified radius grows linearly in the
party count.  It cannot grow faster.

\begin{proposition}[Ceiling on the certified radius]
\label{prop:region-ceiling}
For every \(q\geq2\) and every \(n\) there are traceless Hermitian \(h_j\)
on \(\mathbb C^q\) with
\(\sum_j\lVert h_j\rVert_{\mathrm{op}}=2\pi(1-1/q)n\) and \(D>0\) for which
\(\bigotimes_je^{ih_j}\) is a scalar, hence \(\varepsilon=0\).  No estimate
\(\varepsilon(U)^2\geq\rho D^2/q\) with \(\rho>0\) can therefore hold on a
region of radius exceeding \(2\pi(1-1/q)n\).
\end{proposition}

\begin{proof}
Take \(h_j=(2\pi/q)\diag(q-1,-1,\ldots,-1)\) at every site.  It is
traceless, and both eigenvalues of \(e^{ih_j}\) equal \(e^{-2\pi i/q}\), so
\(e^{ih_j}=e^{-2\pi i/q}I\) and the product of \(n\) of them is a global
phase.  Here \(\lVert h_j\rVert_{\mathrm{op}}=2\pi(q-1)/q\) and
\(\lVert h_j\rVert_F^2=(2\pi/q)^2q(q-1)>0\).
\end{proof}

Linear growth in the party count is therefore the maximal possible order,
and Corollary~\ref{cor:k-uniform-region} attains it for AME states, within
a factor \(4\pi e(1-1/q)\) of the ceiling.  Bounded uniformity order is a
different regime, and the following family measures it exactly.

\begin{proposition}[The certified radius at uniformity order three]
\label{prop:stability-region}
Let \(\ell\geq3\), \(N=2^\ell\), and let \(\lvert C\rangle\) be the
equal-phase CSS state of \(C=\mathrm{RM}(1,\ell)\leq\F_2^N\).  Its
uniformity order is \(3\), for every \(\ell\).  Take
\(h_j=cZ(1)\) at every site and put \(\theta=Nc\in(0,2\pi)\).  Then
\[
 \sum_j\lVert h_j\rVert_{\mathrm{op}}=\theta,\qquad
 D=\frac{\theta\sqrt2}{\sqrt N},\qquad
 \varepsilon=\frac{2\lvert\sin(\theta/2)\rvert}{\sqrt N},\qquad
 \frac D{\sqrt2\,\varepsilon}=\frac\theta{2\sin(\theta/2)} .
\]
The last ratio does not depend on \(N\) and increases on \((0,2\pi)\); it
passes \(\sqrt{6/5}\) at the unique \(\theta^*=1.46562\ldots\).  Hence the
conclusion of Theorem~\ref{thm:two-uniform-stability} fails at
\(t=\theta^*\), for every length.  At uniformity order three the exact
radius for the constant \(\sqrt{6q/5}\) therefore lies between
\(R_3=1/\sqrt2\) and \(\theta^*\).
\end{proposition}

\begin{proof}
The nonzero weights of \(\mathrm{RM}(1,\ell)\) are \(N/2\) and \(N\), and
\(C^\perp=\mathrm{RM}(\ell-2,\ell)\) has minimum distance \(4\).  A
marginal of the equal-phase CSS state fails to be maximally mixed only when
a nonzero word of \(C\) or of \(C^\perp\) is supported inside it, so by
\((2.2)\) the uniformity order is
\(\min(N/2,4)-1=3\) whenever \(N\geq8\).  Since
\(Z(1)\lvert x\rangle=(-1)^x\lvert x\rangle\) at a single site,
\(U\lvert x\rangle=e^{ic(N-2\operatorname{wt}(x))}\lvert x\rangle\), and
the weight distribution \(1,\,2N-2,\,1\) on weights \(0,\,N/2,\,N\) gives
\(\langle U\rangle=1-(1-\cos\theta)/N\).  This is nonnegative, so
\(\varepsilon^2=2-2\langle U\rangle=4\sin^2(\theta/2)/N\).  Also
\(\lVert cZ(1)\rVert_F^2=2c^2\) and
\(\lVert cZ(1)\rVert_{\mathrm{op}}=c\), giving the displayed \(D\) and
\(\theta\).  Finally \(\theta/(2\sin(\theta/2))=u/\sin u\) with
\(u=\theta/2\in(0,\pi)\), which increases from \(1\) to \(\infty\); at
\(\theta^*\) the ratio exceeds the constant of
Theorem~\ref{thm:two-uniform-stability} while the defect is
\(2\sin(\theta^*/2)N^{-1/2}\).
\end{proof}

\AMEFigureFive

Two consequences.  The smallness hypothesis of
Theorem~\ref{thm:two-uniform-stability} is load-bearing, and its constant
\(\frac12\) is within a factor of about three of sharp, since the family
leaves the hypothesis at \(\theta=\frac12\) and violates the conclusion at
\(\theta^*\).  And the party count is not what governs the region.  The
family has \(N\) parties and uniformity order three at every length, so
what it bounds is the radius at order three; the defect at which it breaks
the conclusion is \(1.3379\ldots N^{-1/2}\) only because a bounded
\(\ell^1\) radius spread evenly over \(N\) sites has defect image of order
\(N^{-1/2}\).  When the uniformity order grows with the party count, as it
does for an \(\AME(2m,q)\) state, the admissible generator radius grows with
it; the independent constraint \(D\leq\sqrt q/2\) then binds before
\(t\leq R_m\sim m/e\) does.  This enlarges a region already described in
generator coordinates.  It does not imply that every unitary in a
constant-radius defect ball has such coordinates.  Both statements concern
a known state acted on by product unitaries of bounded generator norm;
neither certifies an unknown device.

For stabilizer AME states there is a second route to a stability estimate,
which buys the same conclusion with a different hypothesis set.  It gives
up the global \(\ell^1\) budget on \(\sum_j\lVert h_j\rVert_{\mathrm{op}}\)
entirely, keeps a per-site condition and the global \(\ell^2\) condition,
and pays only in the constant.  The mechanism is a bipartition bound in
which the Weyl amplitudes of the two halves are compared directly, so no
Taylor remainder appears and no budget is needed to control one.

\begin{proposition}[Half-splitting bound]\label{prop:half-splitting}
Let \(\lvert\psi\rangle\) be a stabilizer \(\AME(2m,q)\) state with label
group \(L\) and character \(\chi\), let \(U=\bigotimes_jU_j\) be a product
unitary, and put \(a_j(v)=q^{-1}\Tr(W_v^\dagger U_j)\), so that
\(p_j=\lvert a_j\rvert^2\) is a probability distribution on \(\F_q^2\) for
each \(j\).  For a set \(A\) of \(m\) parties define \(u,v\in\mathbb R^L\)
by
\[
 u(\ell)=\prod_{j\in A}\lvert a_j(\ell_j)\rvert,\qquad
 v(\ell)=\prod_{j\notin A}\lvert a_j(\ell_j)\rvert .
\]
Then \(\lVert u\rVert=\lVert v\rVert=1\) and
\(\lvert\langle\psi\rvert U\lvert\psi\rangle\rvert\leq\langle u,v\rangle\),
so \(\varepsilon(U)^2\geq\lVert u-v\rVert^2\).
\end{proposition}

\begin{proof}
The coordinate restriction \(\pi_A:L\to(\F_q^2)^A\) has kernel the labels
supported on the complementary \(m\)-set, which by \((4.1)\) is trivial;
both sides have \(q^{2m}\) elements, so \(\pi_A\) is a bijection, and
likewise for \(\pi_{A^c}\).  Hence
\(\lVert u\rVert^2=\sum_{x\in(\F_q^2)^A}\prod_{j\in A}p_j(x_j)=1\), and the
same for \(v\).  From \((2.2)\) and
\(\Tr(W_\ell U)=\prod_j\Tr(W_{\ell_j}U_j)\), taking moduli termwise gives
\(\lvert\langle\psi\rvert U\lvert\psi\rangle\rvert
\leq\sum_{\ell\in L}\prod_j\lvert a_j(\ell_j)\rvert=\langle u,v\rangle\),
the relabelling \(\ell\mapsto-\ell\) being a bijection of \(L\).  For unit
vectors \(\langle u,v\rangle=1-\frac12\lVert u-v\rVert^2\).
\end{proof}

The AME condition supplies more than the two bijections: it pins down a
family of labels of minimum weight that straddles the cut, and those are
the labels on which \(u\) and \(v\) are forced apart.

\begin{lemma}[Cut-transversal labels]\label{lem:cut-transversal}
With \(A\) as above and \(j\notin A\), let
\(L_j=\{\ell\in L:\supp\ell\subseteq A\cup\{j\}\}\).  Then
\(\lvert L_j\rvert=q^2\), every nonzero \(\ell\in L_j\) has support exactly
\(A\cup\{j\}\), and for every \(i\in A\cup\{j\}\) the evaluation
\(\ell\mapsto\ell_i\) is a bijection \(L_j\to\F_q^2\).
\end{lemma}

\begin{proof}
\(L_j\) is the kernel of the projection of \(L\) to the coordinates
\(A^c\setminus\{j\}\), which is surjective because it factors the bijection
\(\pi_{A^c}\); so \(\lvert L_j\rvert=q^{2m}/q^{2(m-1)}=q^2\).  A nonzero
\(\ell\in L_j\) has weight at least \(m+1\) by \((4.1)\) and support inside
a set of that size.  For \(i\in A\cup\{j\}\) the kernel of
\(\ell\mapsto\ell_i\) on \(L_j\) consists of labels supported on \(m\)
parties, hence is trivial, and the two sets have the same cardinality.
\end{proof}

\begin{theorem}[Stability without a generator budget]
\label{thm:budget-free-stability}
Let \(\lvert\psi\rangle\) be a stabilizer \(\AME(2m,q)\) state with
\(m\geq3\), let \(U=\bigotimes_jU_j\) be a product unitary, and put
\(\eta_j=1-\lvert q^{-1}\Tr U_j\rvert^2\) and \(H=\sum_j\eta_j\).  If
\(H\leq\frac14\) then \(\varepsilon(U)^2\geq H/4\).  Consequently, if
\(U=\bigotimes_je^{ih_j}\) with each \(h_j\) traceless Hermitian of
spectral spread at most \(\pi\) and \(D\leq\sqrt q/2\), then
\[
 \varepsilon(U)^2\geq\frac{c\,D^2}{4q},
 \qquad\text{that is}\qquad
 D\leq2\sqrt{q/c}\;\varepsilon(U),
\]
with \(c=c(s)\) as in Lemma~\ref{lem:site-contraction}.  No bound on
\(\sum_j\lVert h_j\rVert_{\mathrm{op}}\) is assumed.
\end{theorem}

\begin{proof}
Fix an \(m\)-set \(A\) and write \(\eta=\max_j\eta_j\leq H\).  For
\(j\notin A\) and \(w\in\F_q^2\) let \(\ell^{j,w}\) be the label of
Lemma~\ref{lem:cut-transversal} with \(j\)-coordinate \(w\).  Since
\(\ell^{j,w}\) vanishes off \(A\cup\{j\}\),
\[
 \sum_{w\ne0}v(\ell^{j,w})^2
  =\eta_j\prod_{i\notin A\cup\{j\}}(1-\eta_i)\geq\eta_j(1-H).
\]
Every \(A\)-coordinate of \(\ell^{j,w}\) is nonzero, so
\(p_i(\ell_i^{j,w})\leq\eta_i\) for \(i\in A\); fixing one \(i_0\in A\) and
using injectivity of \(w\mapsto\ell_{i_0}^{j,w}\),
\(\sum_{w\ne0}u(\ell^{j,w})^2\leq\prod_{i\in A}\eta_i\).  Summing over the
\(m\) choices of \(j\) and using
\(\prod_{i\in A}\eta_i\leq\eta^{m-1}m^{-1}\sum_{i\in A}\eta_i\) bounds the
total by \(\eta^{m-1}\sum_{i\in A}\eta_i\).  All terms of
\(\lVert u-v\rVert^2\) are nonnegative, and \((x-y)^2\geq x^2/2-y^2\), so
this family contributes at least
\(\frac12(1-H)\sum_{j\notin A}\eta_j-\eta^{m-1}\sum_{i\in A}\eta_i\).  The
mirror family, supported on \(A^c\cup\{i\}\) for \(i\in A\), is disjoint
from it for \(m\geq2\) and contributes the reflected bound.  Adding,
\[
 \varepsilon(U)^2\geq\lVert u-v\rVert^2
   \geq\Bigl(\tfrac12(1-H)-\eta^{m-1}\Bigr)H .
\]
At \(H\leq\frac14\) and \(m\geq3\) one has
\(\eta^{m-1}\leq H^2\leq\frac1{16}\), so the bracket is at least
\(\frac5{16}\geq\frac14\).  For the second statement,
Lemma~\ref{lem:site-contraction} gives
\(\eta_j\geq c\lVert h_j\rVert_F^2/q\), hence \(H\geq cD^2/q\); and
\(1-\lvert\mathbb E\,e^{iX}\rvert^2\leq\frac12\mathbb E(X-X')^2
=\lVert h_j\rVert_F^2/q\) gives \(H\leq D^2/q\leq\frac14\).
\end{proof}

The two routes are a choice of hypothesis set.
Corollary~\ref{cor:k-uniform-region} assumes per-site spread \(\frac12\),
\(D\leq\sqrt q/2\), and the \(\ell^1\) budget \(t\leq R_k\), and delivers
the constant \(\sqrt{6q/5}\).  Theorem~\ref{thm:budget-free-stability}
keeps \(D\leq\sqrt q/2\), relaxes the spread hypothesis to \(\pi\), deletes
the budget, and pays \(2\sqrt{q/c}\), which is \(\pi\sqrt q\) at spread
\(\pi\) and \(2.02\sqrt q\) at spread \(\frac12\).  Its per-site spread
condition is not decorative: the family of
Proposition~\ref{prop:region-ceiling} has \(e^{ih_j}\) scalar, so every
\(\eta_j\) vanishes and the first bound is empty, and those \(h_j\) have
spread \(2\pi\), exactly where \(c(s)\leq0\) and the second bound is empty
too.  The spread condition is the residue of the ceiling.

One configuration is worth recording, because it is the natural candidate
for a product unitary that Proposition~\ref{prop:half-splitting} might fail
to see.  At \(q=2\) take \(U_j=H\) at every site, so that each \(p_j\) is
uniform on \(\{X,Z\}\); then \(u(\ell)v(\ell)=2^{-m}\) when every
\(\ell_j\in\{X,Z\}\) and vanishes otherwise, so
\(\langle u,v\rangle=2^{-m}\lvert L\cap\{X,Z\}^{2m}\rvert\).  Writing
\(\{X,Z\}=X+\{0,Y\}\), the set \(L\cap\{X,Z\}^{2m}\) is empty or a coset of
\(L\cap\{0,Y\}^{2m}\), and the latter is a binary linear code of length
\(2m\) whose nonzero words have weight at least \(m+1\).  Singleton bounds
its dimension by \(m\), with equality only for a binary
\([2m,m,m+1]\) MDS code; the binary MDS codes are the trivial ones, which
forces \(m\leq1\).  Hence \(\langle u,v\rangle\leq\frac12\) and
\(\varepsilon(H^{\otimes2m})^2\geq1\) for every stabilizer
\(\AME(2m,2)\) state with \(m\geq2\).  The bound is not merely nonvacuous
there; it is within \(\sqrt2\) of the largest possible defect.  What the
half-splitting bound discards is phase information, but the mechanism that
detects this configuration is a count of codewords in a structured subset
of \(L\), and the MDS constraints force that count down.

\begin{corollary}[Approximate symmetries decompose]
\label{cor:approximate-decomposition}
Let \(\lvert\psi\rangle\) be a stabilizer \(\AME(2m,q)\) state with
\(m\geq2\).  There is an \(\varepsilon_0(\psi)>0\) such that every product
unitary \(U\) with \(\varepsilon(U)<\varepsilon_0\) factors as
\[
 U=g\cdot\bigotimes_{j=1}^{2m}e^{ih_j},\qquad g\in G(\psi),
\]
with the \(h_j\) traceless Hermitian,
\(\sum_j\lVert h_j\rVert_{\mathrm{op}}\leq\frac12\), and
\[
 \Bigl(\sum_j\lVert h_j\rVert_F^2\Bigr)^{1/2}
   \leq\sqrt{6q/5}\;\varepsilon(U).
\]
By Theorem~\ref{thm:lu-lc-rigidity} the factor \(g\) is an exact
local-Clifford symmetry.
\end{corollary}

The size constraint on the \(h_j\) is part of the conclusion and cannot be
dropped: the generators of Proposition~\ref{prop:region-ceiling} are
traceless with \(\bigotimes_je^{ih_j}\) scalar, so the defect vanishes
while the generator norm does not.

\begin{proof}
Work on the compact manifold \(X=\mathrm{PU}(q)^{2m}\), on which
\(f(U)=\varepsilon(U)^2\) is well defined because the defect depends on
\(U\) only through \(\lvert\langle\psi\rvert U\lvert\psi\rangle\rvert\).

\emph{Quadratic growth.}  Let \(B\subseteq X\) be the image of
\(\{(h_j)\ \text{traceless}:\sum_j\lVert h_j\rVert_{\mathrm{op}}
\leq\frac12\}\) under \((h_j)\mapsto[\bigotimes_je^{ih_j}]\).  That map is
a diffeomorphism near the origin, of matching dimension
\(2m(q^2-1)\), so \(B\) is a neighbourhood of the identity coset, and
Theorem~\ref{thm:two-uniform-stability} gives
\(f\geq\frac5{6q}\sum_j\lVert h_j\rVert_F^2\) on it.  For \(g\in G(\psi)\)
one has \(f(gV)=f(V)\): if \(g\lvert\psi\rangle=e^{i\theta}\lvert\psi\rangle\)
then \(\langle\psi\rvert g V\lvert\psi\rangle
=e^{-i\theta}\langle\psi\rvert V\lvert\psi\rangle\) after conjugating the
left factor, and \(f\) depends only on the modulus.  Hence the same estimate
holds on the left translate \(gB\), with the generators measured from \(g\),
which is the form the factorization \(U=g\cdot\bigotimes_je^{ih_j}\) requires.

\emph{Isolation and the threshold.}  The zero set of \(f\) on \(X\) is the
image \(Z\) of \(G(\psi)\), finite by
Theorem~\ref{thm:two-uniform-discrete} and
Remark~\ref{rem:local-phase-torus}.  So \(\bigcup_{g\in Z}gB\) is a
neighbourhood of \(Z\), its complement \(K\) is compact, and \(f>0\) on
\(K\); put \(\varepsilon_0=(\min_Kf)^{1/2}>0\).  If
\(\varepsilon(U)<\varepsilon_0\) then \(U\notin K\), so \(U\in gB\) for
some \(g\in Z\), and the first step applies at that \(g\).
\end{proof}

The stabilizer hypothesis enters this corollary only in the last sentence
of its statement.  The factorization itself holds for every \(2\)-uniform
state, with \(g\) an exact product-unitary symmetry; and finiteness of the
zero set is a consequence of \((B.1)\) rather than a separate input, since
quadratic growth already isolates the zeros of \(f\) on a compact manifold.
The threshold this proof produces comes from compactness and is not
exhibited; Theorem~\ref{thm:explicit-threshold} below replaces it by a
formula, at the cost of the stabilizer hypothesis and of a factor
exponential in the party count.

Corollary~\ref{cor:approximate-decomposition} measures how far an
approximate symmetry of one state is from an exact one.  The intertwiner
question between two states is different, and its answer is the tool the
explicit threshold needs: it converts a small global defect into a nearby
Clifford at each party, with a threshold exhibited in closed form.

\begin{lemma}[Quantitative diagonal-tensor axes]
\label{lem:quantitative-axes}
Let \(r\geq3\), let \(E_1,\ldots,E_r\) be \(N\)-dimensional complex inner
product spaces with orthonormal bases \(\{e_{ij}\}\), and put
\[
 T=a\sum_{j=1}^N\lambda_je_{1j}\otimes\cdots\otimes e_{rj},
 \qquad a>0,\quad\lvert\lambda_j\rvert=1 .
\]
Let \(T'\) have the same form in orthonormal bases \(\{e'_{ij}\}\) of
spaces \(E'_i\), and let \(A_i:E_i\to E'_i\) be unitaries with
\[
 \lVert(A_1\otimes\cdots\otimes A_r)T-T'\rVert\leq\eta,
 \qquad\delta=\eta/a<1/\sqrt2 .
\]
Then in each factor there is a unique permutation \(\pi_i\) with
\[
 \min_{\lvert\alpha\rvert=1}
 \lVert A_ie_{ij}-\alpha e'_{i,\pi_i(j)}\rVert\leq\sqrt2\,\delta
 \qquad(1\leq j\leq N).                                  \tag{B.4}
\]
\end{lemma}

\begin{proof}
Contract both tensors in the first factor against a target frame vector
\(e_{1k}'^{\,*}\).  The contraction of \(T'\) has rank one, and the
contraction of \((A_1\otimes\cdots\otimes A_r)T\) is within \(\eta\) of it.
Flattening the latter between its first remaining factor and the last
\(r-2\) exhibits its singular values as
\(a\lvert\langle e'_{1k},A_1e_{1j}\rangle\rvert\), \(1\leq j\leq N\).  By
the Eckart--Young theorem the squared mass outside the largest is at most
\(\eta^2\), hence at most \(\delta^2\) after dividing by \(a\).
Unitarity of \(A_1\) and completeness of the source basis make the total
squared mass one, so the dominant coefficient has square at least
\(1-\delta^2>\frac12\).  Two orthogonal rows cannot select the same
dominant column, and the overlap matrix is square, so the selections form a
permutation \(\pi_1\).  Each matched column also has total squared mass one
and contains that coefficient, so its off-axis mass is at most
\(\delta^2\), and \(2\bigl(1-\sqrt{1-\delta^2}\bigr)\leq2\delta^2\) gives
\((B.4)\).  Repeat in every factor.
\end{proof}

\begin{proposition}[Quantitative intertwiner rigidity]
\label{prop:quantitative-intertwiner}
Let \(q=p^e\), let \(m\geq2\), and let
\(\lvert\psi\rangle,\lvert\phi\rangle\) be stabilizer \(\AME(2m,q)\)
states.  Put
\[
 \tau_p=\min\Bigl\{\tfrac13,\ \frac{\sin(\pi/p)}{2\sqrt2}\Bigr\},
 \qquad
 \delta=2q^{(m+1)/2}\varepsilon .
\]
Suppose that, after an allowed party permutation,
\[
 \min_{\lvert\alpha\rvert=1}
 \bigl\lVert(U_1\otimes\cdots\otimes U_{2m})\lvert\psi\rangle
   -\alpha\lvert\phi\rangle\bigr\rVert\leq\varepsilon
\]
and that \(\delta<\tau_p\).  Then for every party \(i\) there is a
single-qudit Clifford \(K_i\) with
\[
 q^{-1/2}\lVert U_i-K_i\rVert_{\mathrm{HS}}
  \leq\sqrt2\,\delta=2\sqrt2\,q^{(m+1)/2}\varepsilon .    \tag{B.5}
\]
For \(m=3\) and equal-phase CSS states of linear \([6,3,4]_q\) MDS codes
this is the bound \(2\sqrt2\,q^2\varepsilon\) under \(2q^2\varepsilon<\tau_p\).
\end{proposition}

\begin{proof}
Fix an \((m+1)\)-party set \(S\).  By \((2.2)\) and
Proposition~\ref{prop:stabilizer-ame-support} the supported label group
\(L_\psi(S)\) has \(q^2\) elements and projects bijectively onto
\(\F_q^2\) at every party of \(S\), so the nonidentity part of the reduced
operator is
\[
 T_S=q^{-(m+1)}\sum_{v\ne0}\lambda_v
   W_{1,L_1v}\otimes\cdots\otimes W_{m+1,L_{m+1}v},
 \qquad\lvert\lambda_v\rvert=1,
\]
where \(v\) runs over the nonzero labels at one fixed party of \(S\) and
each \(L_i\) is the resulting parametrization of \(L_\psi(S)\) at party
\(i\).  Normalizing every Weyl operator by \(q^{-1/2}\) puts \(T_S\) in the
form of Lemma~\ref{lem:quantitative-axes} with
\(r=m+1\geq3\), \(N=q^2-1\), and
\(a=q^{-(m+1)}q^{(m+1)/2}=q^{-(m+1)/2}\).  Two
unit vectors at phase-optimized distance \(\varepsilon\) have density
matrices at trace distance at most \(2\varepsilon\); partial trace
contracts the trace norm, and the Hilbert--Schmidt norm is no larger, so
every \((m+1)\)-party marginal satisfies \(\eta\leq2\varepsilon\) and hence
\(\eta/a\leq2q^{(m+1)/2}\varepsilon=\delta\).  The lemma then matches Weyl
axes within \(\kappa=\sqrt2\,\delta\).

It remains to turn approximate axes into an exact Clifford, party by party.
Extend the matched label permutation by \(\pi(0)=0\); the identity axis is
matched exactly.  Conjugating a product of two Weyl operators in the two
available orders places the axes labelled \(\pi(v+w)\) and
\(\pi(v)+\pi(w)\) within \(3\kappa\).  Distinct Weyl axes are at
phase-minimized distance \(\sqrt2\), so \(\delta<\frac13\) forces
\(\pi(v+w)=\pi(v)+\pi(w)\); the matching was a permutation, so \(\pi\) is
\(\F_p\)-linear.  Conjugation preserves Weyl commutators exactly, replacing
the four factors of a commutator by matched axes costs at most
\(4\kappa\), and distinct \(p\)-th root commutator phases are separated by
\(2\sin(\pi/p)\), so \(\delta<\sin(\pi/p)/(2\sqrt2)\) forces \(\pi\) to
preserve the prime-field symplectic form.  By \((2.0)\) it is realized by
an exact Clifford \(K_i\).

Finally the phases.  Put \(V=K_i^\dagger U_i\) and expand
\(V=\sum_xc_xW_x\) with \(\sum_x\lvert c_x\rvert^2=1\).  Axis closeness
gives \(\lvert q^{-1}\Tr(W_v^\dagger VW_vV^\dagger)\rvert\geq
1-\kappa^2/2\) for every \(v\), and character orthogonality turns this into
\(\sum_x\lvert c_x\rvert^4\geq(1-\kappa^2/2)^2\).  Hence some
\(\lvert c_x\rvert\geq1-\kappa^2/2\), so \(V\) lies within \(\kappa\) of a
phase times \(W_x\); absorbing that displacement into \(K_i\) gives
\((B.5)\).  Every party lies in some \((m+1)\)-set, as in
Corollary~\ref{cor:full-weyl-cover}.
\end{proof}

The exponent in \(\delta\) is a feature of the states rather than slack in
the estimate.  The nonidentity part of an \((m+1)\)-party marginal of an
\(\AME(2m,q)\) state has Hilbert--Schmidt norm
\(q^{-(m+1)/2}\sqrt{q^2-1}\): the marginal is nearly maximally mixed, and
the Weyl signal that exposes the local factors is that small, so a global
defect is diluted to that scale before any marginal can read it off.

The closest general framework is the perturbation theory of nearly
orthogonally decomposable tensors, which gives dimension-independent bounds
without a singular-value gap \cite{AuddyYuan2020}; the argument above is
the equal-weight complex specialization needed here, together with the
Weyl multiplication, commutator, and implementing-Clifford steps.  Its
conclusion is distance from the ambient Clifford group and nothing finer.
Whenever two states admit an exact intertwiner outside a proposed smaller
subgroup, that intertwiner has positive distance from the finite closed
subgroup, so taking \(\varepsilon=0\) falsifies every uniform estimate
toward it.

Rounding a product unitary to a nearby product Clifford is only half of a
decomposition.  The other half is deciding whether that Clifford is an
exact symmetry, and for stabilizer states it needs no approximation
argument at all, because their overlaps take only finitely many values.

\begin{lemma}[Quantized stabilizer overlap]
\label{lem:stabilizer-overlap-gap}
Let \(\lvert\psi\rangle,\lvert\phi\rangle\) be stabilizer states on \(n\)
qudits of local dimension \(q=p^e\), with label groups \(L_\psi,L_\phi\)
and stabilizer characters \(\chi_\psi,\chi_\phi\).  Then
\[
 \lvert\langle\psi\vert\phi\rangle\rvert^2=
 \begin{cases}
  \lvert L_\psi\cap L_\phi\rvert/q^n,
   &\text{if \(\chi_\psi=\chi_\phi\) on \(L_\psi\cap L_\phi\)},\\
  0,&\text{otherwise.}
 \end{cases}
\]
In particular
\(\lvert\langle\psi\vert\phi\rangle\rvert^2\in\{0\}\cup\{p^{-j}:j\geq0\}\),
so \(\lvert\langle\psi\vert\phi\rangle\rvert>p^{-1/2}\) forces
\(\lvert\phi\rangle\in\mathbb C\lvert\psi\rangle\).
\end{lemma}

\begin{proof}
By \((2.2)\), \(\rho_\psi=q^{-n}\sum_{\ell\in L_\psi}\chi_\psi(\ell)W_\ell\)
and likewise for \(\phi\).  Weyl orthogonality
\(\Tr(W_\ell^\dagger W_{\ell'})=q^n\delta_{\ell,\ell'}\) gives
\[
 \lvert\langle\psi\vert\phi\rangle\rvert^2=\Tr(\rho_\psi\rho_\phi)
  =q^{-n}\sum_{\ell\in L_\psi\cap L_\phi}
    (\chi_\psi\overline{\chi_\phi})(\ell).
\]
The summand is a character of the group \(L_\psi\cap L_\phi\), so the sum
is \(\lvert L_\psi\cap L_\phi\rvert\) when that character is trivial and
\(0\) otherwise.  The group has \(p\)-power order and \(q^n=p^{en}\), which
gives the quantization; the value \(1\) occurs only when \(L_\psi=L_\phi\)
and the characters agree, that is, when the two states span the same ray.
\end{proof}

The quantization has a consequence about the symmetry group that is
stronger, and easier to remember, than any threshold formula.
Theorem~\ref{thm:two-uniform-discrete} makes \(G(\psi)\) finite, hence its
elements isolated; on the Clifford locus they are isolated by a fixed
amount that no parameter of the problem can erode.

\begin{corollary}[Uniform separation of the exact symmetries]
\label{cor:uniform-separation}
Let \(\lvert\psi\rangle\) be a stabilizer state on \(n\) qudits of local
dimension \(q=p^e\), and let \(K\) be a product Clifford.  Then either
\(K\in G(\psi)\) or
\[
 \varepsilon(K)\geq\bigl(2-2p^{-1/2}\bigr)^{1/2}\geq\bigl(2-\sqrt2\bigr)^{1/2}
  =0.7653\ldots
\]
Equivalently, if \(K\) and \(K'\) are product Cliffords in different cosets
of \(G(\psi)\), then \(K^\dagger K'\) has defect at least this much.  The
bound depends on the characteristic alone: not on the number of parties,
not on the state, and not on the local dimension beyond its
characteristic.
\end{corollary}

\begin{proof}
\(K\lvert\psi\rangle\) is a stabilizer state, so
Lemma~\ref{lem:stabilizer-overlap-gap} applies to the pair
\(\lvert\psi\rangle,K\lvert\psi\rangle\): either
\(\lvert\langle\psi\rvert K\lvert\psi\rangle\rvert\leq p^{-1/2}\), whence
\(\varepsilon(K)^2=2-2\lvert\langle\psi\rvert K\lvert\psi\rangle\rvert
\geq2-2p^{-1/2}\), or the overlap exceeds \(p^{-1/2}\) and
\(K\lvert\psi\rangle\in\mathbb C\lvert\psi\rangle\).  For two product
Cliffords apply this to \(K^\dagger K'\), which is again one.  Since
\(p\geq2\), the bound is least at \(p=2\).
\end{proof}

So on the Clifford locus the defect has no small nonzero values.  This is
what removes the compactness step from
Corollary~\ref{cor:approximate-decomposition}: a product unitary with small
defect is first rounded to a product Clifford by
Proposition~\ref{prop:quantitative-intertwiner}, and the rounded operator
is then an exact symmetry outright, because a near-miss is impossible.

The rounding need not inspect one minimum-support marginal in isolation.
All such marginals are projections of orthogonal exact-support sectors of
the same global density-matrix difference.  Summing their squared errors
therefore recovers part of the signal lost by marginalization.

\begin{lemma}[Collective minimum-support energy]
\label{lem:collective-support-energy}
Let \(\rho=\lvert\psi\rangle\langle\psi\rvert\),
\(\Delta=U\rho U^\dagger-\rho\), \(n=2m\), and \(r=m+1\).  Then
\[
 \sum_{\substack{A\subseteq[n]\\|A|=r}}
 q^r\lVert\Delta_A\rVert_{\mathrm{HS}}^2
 \leq q^n\lVert\Delta\rVert_{\mathrm{HS}}^2
 \leq2q^n\varepsilon(U)^2.                       \tag{B.7}
\]
Writing \(\delta_A=q^{r/2}\lVert\Delta_A\rVert_{\mathrm{HS}}\) and
\(N_m=\binom{2m}{m+1}\), there are two \(r\)-sets \(A,B\) with
\(A\cup B=[n]\) and
\[
 \delta_A^2+\delta_B^2
 \leq \frac{4q^n}{N_m}\varepsilon(U)^2.          \tag{B.8}
\]
\end{lemma}

\begin{proof}
The stabilizer expansion of \(\rho\) has no nonidentity component on at
most \(m\) sites.  Local conjugation preserves exact support because it
preserves tracelessness.  Thus \(\Delta_A\), for \(|A|=r\), is precisely
the partial trace of the exact-support-\(A\) component of \(\Delta\).
If that component is \(D_A\otimes I_{A^c}\), then
\(\Delta_A=q^{n-r}D_A\) and
\[
 \lVert D_A\otimes I_{A^c}\rVert_{mathrm{HS}}^2
 =q^{-(n-r)}\lVert\Delta_A\rVert_{mathrm{HS}}^2.
\]
Exact-support subspaces are mutually orthogonal, which gives the first
inequality in (B.7).  The second follows from
\(\lVert\Delta\rVert_{mathrm{HS}}^2
=2(1-|\langle\psi|U|\psi\rangle|^2)leq2\varepsilon(U)^2\).

Join two \(r\)-sets when their union is \([n]\).  This graph is regular:
for fixed \(A\), the complement of a neighbor is any \((m-1)\)-subset of
\(A\).  Averaging \(\delta_A^2+\delta_B^2\) over its edges gives twice
the vertex average.  One edge therefore satisfies (B.8).
\end{proof}

The two good sets cover every party, so the quantitative-axis argument can
be applied to \(A\) and \(B\), using \(A\) on their intersection.  This
gives one product Clifford while retaining collective \(\ell^2\) control
of the local errors.  Two-uniformity then converts that control directly
to state distance, without a telescoping \(\ell^1\) estimate.

\begin{theorem}[Aggregate global rounding]
\label{thm:aggregate-global-rounding}
Let \(\lvert\psi\rangle\) be a stabilizer \(\AME(2m,q)\) state with
\(q=p^e\) and \(m\geq3\).  Put
\[
 \begin{gathered}
 N_m=\binom{2m}{m+1},\quad r=m+1,\quad
 d_p=(2-2p^{-1/2})^{1/2},\\
 R_{m,q}=\min\left\{
 \frac{\tau_p\sqrt{N_m}}{2q^m},
 \frac{\sqrt{N_m}}{2q^{m+1/2}},
 \frac{\sqrt{N_m}}{2\pi\sqrt{2r}\,q^m},
 \frac{d_p}{1+\pi\sqrt{2r q^{2m}/N_m}}
 \right\}.                                             \tag{B.9}
 \end{gathered}
\]
If \(\varepsilon(U)<R_{m,q}\), then, up to a global phase,
\[
 U=g\bigotimes_{i=1}^{2m}e^{ih_i},\qquad g\in G(\psi),
\]
where the \(h_i\) are traceless Hermitian, have spectral spread at most
\(\pi\), and satisfy
\[
 \left(\sum_i\lVert h_i\rVert_F^2\right)^{1/2}
 \leq\pi\sqrt q\,\varepsilon(U).                       \tag{B.10}
\]
Moreover \(R_{m,q}=\Theta_q(m^{-3/4}(2/q)^m)\).  Thus this route is
polynomial in \(m\) at \(q=2\), whenever the stated AME state exists,
improves the exponential base at \(q=3\),
improves the frame-rounding prefactor at \(q=4\), and is superseded by
Theorem~\ref{thm:explicit-threshold} for \(q>4\).
\end{theorem}

\begin{proof}
Choose \(A,B\) from Lemma~\ref{lem:collective-support-energy}.  The
nonidentity part of each marginal has diagonal Weyl coefficient
\(q^{-r/2}\), so Lemma~\ref{lem:quantitative-axes} and its
additivity, commutator, and phase steps give a Clifford \(K_i\) at each
party with
\[
 q^{-1/2}\lVert U_i-K_i\rVert_{\mathrm{HS}}
 \leq\sqrt2\delta_i,\qquad
 \sum_i\delta_i^2
 \leq r(\delta_A^2+\delta_B^2)
 \leq\frac{4rq^n}{N_m}\varepsilon(U)^2,          \tag{B.11}
\]
where \(\delta_i=\delta_A\) on \(A\) and \(\delta_i=\delta_B\) on
\(A^c\subset B\).  The first clause of (B.9) supplies the axis threshold
\(\tau_p\).

Choose principal eigenangles of \(K_i^\dagger U_i\) after optimizing its
scalar phase and subtract their mean, obtaining traceless \(h_i\).
Since \(|x|\leq(\pi/2)|e^{ix}-1|\) for \(|x|\leq\pi\),
\[
 D^2:=\sum_i\lVert h_i\rVert_F^2
 \leq\frac{\pi^2q}{2}\sum_i\delta_i^2
 \leq\frac{2\pi^2qrq^n}{N_m}\varepsilon(U)^2.   \tag{B.12}
\]
The second clause of (B.9) puts all eigenangles in an arc of length
\(\pi\); the third gives \(D\leq\sqrt q/2\).

Put \(K=\bigotimes_iK_i\) and
\(V=K^\dagger U=\bigotimes_ie^{ih_i}\), up to phase.  If
\(M=\sum_ih_i^{(i)}\), two-uniformity and tracelessness give
\(\langle M^2\rangle=D^2/q\).  Hence
\[
 \lVert(V-I)\psi\rVert^2
 =\langle2-2\cos M\rangle\leq\langle M^2\rangle=D^2/q.
\]
Together with (B.12), the fourth clause of (B.9) puts the rays of
\(K\psi\) and \(\psi\) less than \(d_p\) apart.  Corollary
\ref{cor:uniform-separation} forces \(K\in G(\psi)\).  Apply
Theorem~\ref{thm:budget-free-stability} to the residual \(V\): its
constant has \(c\geq4/\pi^2\), giving (B.10).

Finally \(N_m\sim4^m/\sqrt{\pi m}\).  The third and fourth clauses of
(B.9) have order \(m^{-3/4}(2/q)^m\), and the other two are larger by a
polynomial factor.  Comparison with the
\(q^{-(m+2)/2}/m\) scale below gives the stated boundary at \(q=4\).
\end{proof}

The theorem separates frame recognition from the choice of residual
metric.  If the \(\ell^1\) conclusion of
Corollary~\ref{cor:approximate-decomposition} is required, the same
covering pair and a telescoping estimate give it below
\[
 \min\left\{
  \frac{\tau_p\sqrt{N_m}}{2q^m},
  \frac{\sqrt{N_m}}
       {4\pi\sqrt{2q}\sqrt{(m+1)^2+(m-1)^2}\,q^m}
 \right\}.                                           \tag{B.13}
\]
At \(q=2\) this is of order \(m^{-5/4}\).  The loss relative to (B.9)
comes from imposing a global generator \(\ell^1\) budget after the frames
are known; it is not a remaining marginal-dilution loss.

\subsection{The single-marginal decomposition threshold}

\begin{theorem}[Explicit decomposition threshold]
\label{thm:explicit-threshold}
Let \(\lvert\psi\rangle\) be a stabilizer \(\AME(2m,q)\) state with
\(q=p^e\), \(m\geq2\), \(n=2m\).  Put
\[
 \varepsilon_0=\min\Bigl\{
   \frac{\tau_p}{2q^{(m+1)/2}},\
   \frac1{4\sqrt2\,\pi n\,q^{(m+2)/2}}\Bigr\},
 \qquad
 \tau_p=\min\Bigl\{\tfrac13,\frac{\sin(\pi/p)}{2\sqrt2}\Bigr\}.
\]
Every product unitary \(U\) with \(\varepsilon(U)<\varepsilon_0\) satisfies
the conclusion of Corollary~\ref{cor:approximate-decomposition}.
\end{theorem}

\begin{proof}
Write \(\varepsilon=\varepsilon(U)\) and
\(B=2\sqrt2\,n\,q^{(m+2)/2}\varepsilon\), so that
\(\varepsilon<\varepsilon_0\) gives \(B<1/(2\pi)\).

\emph{Rounding.}  Since \(2q^{(m+1)/2}\varepsilon<\tau_p\),
Proposition~\ref{prop:quantitative-intertwiner} applied with
\(\lvert\phi\rangle=\lvert\psi\rangle\) and the identity permutation
supplies Cliffords \(K_i\) with
\(\lVert U_i-K_i\rVert_{\mathrm{op}}\leq
\lVert U_i-K_i\rVert_{\mathrm{HS}}\leq2\sqrt2\,q^{(m+2)/2}\varepsilon\).
Summing over the \(n\) parties gives
\(\sum_i\lVert U_i-K_i\rVert_{\mathrm{op}}\leq B\).

\emph{The rounded Clifford is exact.}  Put \(K=\bigotimes_iK_i\), so that
\(\lVert U-K\rVert_{\mathrm{op}}\leq B\) and
\(\lvert\langle\psi\rvert K\lvert\psi\rangle\rvert\geq
1-\varepsilon^2/2-B\).  Here \(B<1/(2\pi)<0.1592\), and
\(\varepsilon\leq\tau_p/(2q^{(m+1)/2})\leq1/(6\cdot2^{3/2})<0.06\) because
\(q\geq2\) and \(m\geq2\), so the right side exceeds
\(1-0.0018-0.1592>1/\sqrt2\geq p^{-1/2}\).  By
Corollary~\ref{cor:uniform-separation}, \(K\in G(\psi)\); set \(g=K\).

\emph{The residual generators are small.}  Each factor of
\(V=g^\dagger U\) satisfies
\(\min_\varphi\lVert K_i^\dagger U_i-e^{i\varphi}I\rVert_{\mathrm{op}}
\leq\lVert U_i-K_i\rVert_{\mathrm{op}}\leq B<1\).  For a unitary \(W\) on
\(\mathbb C^q\) with
\(\sigma=\min_\varphi\lVert W-e^{i\varphi}I\rVert_{\mathrm{op}}\leq1\), the
eigenvalues of \(e^{-i\varphi}W\) at the optimal phase are \(e^{i\mu_a}\)
with \(2\lvert\sin(\mu_a/2)\rvert\leq\sigma\) and
\(\mu_a\in(-\pi,\pi]\), so
\(\lvert\mu_a\rvert\leq2\arcsin(\sigma/2)\leq\pi\sigma/2\); subtracting the
mean gives \(W=e^{i\varphi'}e^{ih}\) with \(h\) traceless Hermitian and
\(\lVert h\rVert_{\mathrm{op}}\leq2\max_a\lvert\mu_a\rvert\leq\pi\sigma\).
Applying this at each site,
\(\sum_i\lVert h_i\rVert_{\mathrm{op}}\leq\pi B<\frac12\), and the
collected phases are a global phase absorbed into \(g\).  Finally
\(\varepsilon(V)=\varepsilon(U)\), because \(g\) fixes the ray of
\(\lvert\psi\rangle\), so Theorem~\ref{thm:two-uniform-stability} applies
to \(V\) on its hypothesis and gives
\(D\leq\sqrt{6q/5}\,\varepsilon(U)\).
\end{proof}

Nothing in this proof is a compactness argument: the finite zero set is
never enumerated and no minimum over a compact remainder is taken.  What
compactness supplied in Corollary~\ref{cor:approximate-decomposition} is
supplied here by Corollary~\ref{cor:uniform-separation}, in closed form.
The factorization is also constructive in the sense that matters: \(g\) is
produced from \(U\) by rounding each local factor to a Clifford, and an
\((m+1)\)-party marginal certifies the rounding.

For this single-marginal route the price is the size of \(\varepsilon_0\), which is
\(q^{-(n+4)/4}/n\) up to constants and therefore exponentially small in the
number of parties.  Two steps pay it.  Rounding costs \(q^{(m+1)/2}\), and
by the marginal-dilution remark after
Proposition~\ref{prop:quantitative-intertwiner} no detector that reads a
single small marginal can avoid that factor.  Summing \(n\) per-site
operator-norm errors and taking \(n\) logarithms costs a further
\(n\sqrt q\), which is the price of a per-site conclusion.

The exponential gap left by this theorem and
Theorem~\ref{thm:aggregate-global-rounding} is closed by
Theorem~\ref{thm:cleaning-global-rounding}.  The remaining comparison is
with Proposition~\ref{prop:stability-region}.  That proposition bounds
any threshold valid across states of uniformity order three; an
\(\AME(2m,q)\) state has uniformity order \(m\), and for such states
Corollary~\ref{cor:k-uniform-region} certifies the quadratic estimate on a
generator-coordinate region whose \(\ell^1\) radius grows with \(m\), as
does Theorem~\ref{thm:budget-free-stability} once \(m\geq3\).  Those
estimates control defect from specified generators; they do not put a
defect ball inside that coordinate region.  The cleaning theorem supplies
the inverse statement at radius of order \(n^{-1/2}\), with no claim that
this party-count dependence is optimal.

\subsection{Local gauge groups of two-unitary gates}

\begin{corollary}[Local gauge groups of \(2\)-unitary gates]
\label{cor:two-unitary-gauge}
Let \(W\) be a two-qudit gate whose associated four-party state
\(\psi_{ijkl}=q^{-1}W_{kl,ij}\) is \(\AME(4,q)\) --- equivalently, \(W\) is
\(2\)-unitary, remaining unitary after realignment and after partial
transposition --- and let \(G(\psi)\) be its product-unitary symmetry
group.  Then
\((u_1,u_2,u_3,u_4)\mapsto(\overline u_1,\overline u_2,u_3,u_4)\) is an
isometric group isomorphism carrying \(G(\psi)\) onto the local gauge
group
\[
 \{(v_1\otimes v_2,\,u_3\otimes u_4):
   (u_3\otimes u_4)W(v_1\otimes v_2)^\dagger=e^{i\varphi}W\},
\]
and it preserves \(\bigl(\sum_j\lVert h_j\rVert_F^2\bigr)^{1/2}\).
Consequently the local gauge group of any \(2\)-unitary gate is finite
modulo the local phases, Theorem~\ref{thm:two-uniform-stability} transports
with its hypothesis \(\sum_j\lVert h_j\rVert_{\mathrm{op}}\leq\frac12\),
and Corollary~\ref{cor:approximate-decomposition} transports with its
threshold; for a gate arising from a stabilizer \(\AME(4,q)\) state
Theorem~\ref{thm:explicit-threshold} exhibits that threshold at \(m=2\),
namely \(\varepsilon_0=\min\{\tau_p/(2q^{3/2}),\,
1/(16\sqrt2\,\pi q^2)\}\).  Gate defects
are measured in the metric \(q^{-1}\lVert\cdot\rVert_F\).  For a gate
arising from a stabilizer \(\AME(4,q)\) state the exact gauge group is
local Clifford.
\end{corollary}

\begin{proof}
Normalization is consistent: \(\lVert W\rVert_F^2=q^2\) for a unitary on
\(\mathbb C^q\otimes\mathbb C^q\), so
\(\langle\psi\vert\psi\rangle=\lVert W\rVert_F^2/q^2=1\).  A product
unitary acts on the reshaped gate by
\(W\mapsto(u_3\otimes u_4)W(u_1\otimes u_2)^{T}\), so
\(\lVert U\lvert\psi\rangle-e^{i\theta}\lvert\psi\rangle\rVert
 =q^{-1}\lVert(u_3\otimes u_4)W(u_1\otimes u_2)^T-e^{i\theta}W\rVert_F\);
this is the stated metric.  Transposition preserves the Frobenius norm and
\(u\mapsto\overline u\) is an automorphism of \(U(q)\) fixing
\(\lVert h\rVert_F\) for the corresponding generators, and
\((u_1\otimes u_2)^T=(v_1\otimes v_2)^\dagger\) with
\(v_a=\overline u_a\), which gives the displayed isomorphism.  An
\(\AME(4,q)\) state is \(2\)-uniform, so
Theorem~\ref{thm:two-uniform-discrete} applies, with
Remark~\ref{rem:local-phase-torus} identifying the quotient as the one by
the local phases; the stabilizer case is
Theorem~\ref{thm:lu-lc-rigidity}.
\end{proof}

The statement is vacuous at \(q=2\), where no \(\AME(4,2)\) state exists.
In the dual-unitary circuit literature the local gauge freedom is the
standing nuisance parameter of the classification programme
\cite{RatherAravindaLakshminarayan2022}, and the finite-field Clifford
subfamily has been classified, with perfect-tensor cores appearing among
its equivalence classes \cite{Pahari2026}; the corollary says that for
\(2\)-unitary gates the gauge group itself is finite and, in the stabilizer
case, Clifford.
